\section{An optional precision family for the scalar tree}

The mixed-weight scalar construction has a simple general form. This
lemma describes a precision--coordinate tradeoff within this particular
tree; it makes no optimality claim about scalar certificates in general.

Retain the rational functions $A,B,C,D_0,T$ and the positive scalar
contact product $r$ from the preceding scalar proof. Assume the retained
point and proxy domains, including the nonzero shifted-point guard.
For an integer $n\ge1$, put
\[
 u=2^{n-1},\qquad M=2^{n+1}-1=4u-1,\qquad
 \kappa=r^{-M},\qquad y=T/r^{M+1}.
\]
Keep
\[
 n_0=1+A^2/r,\quad n_1=B^2/r^2+C^2/r,\quad
 L=n_0^2+r n_1^2,
\]
and set
\[
 J=D_0^2+\kappa y^2,\qquad W=L^2+\kappa J^2,
 \qquad b=42rW^2.
\]
Starting from $q_1=r^{-1}$, the $n$ quadratic Laurent equations
\[
 rq_{2^{j+1}-1}=q_{2^j-1}^{\,2}\quad(1\le j\le n)
\]
give $q_M=r^{-M}$; the single quadratic equation $ry=Tq_M$ gives
the required scaled proxy. The norm nodes have the same quadratic
realizations as in the selected construction.

\paragraph{Uniform scalar order.}
At a consumed prime let $R=v_p(r)>0$ be odd. The pole predicate gives
$v_p(W)>-R/2$. Both $R$ and $-MR$ are odd, so the two summands in
each norm node have unequal valuation parity. Thus
\[
 v_p(L)>-R/4,\qquad
 v_p(J)>\frac{2M-1}{4}R.
\]
The unchanged left subtree gives
\[
 v_p(A)>7R/16,\qquad v_p(B)>11R/16,\qquad v_p(C)>3R/16.
\]
The right subtree gives
\[
 v_p(D_0)>\frac{2M-1}{8}R>0,\qquad
 v_p(y)>\frac{6M-1}{8}R.
\]
All first four scalar coefficients therefore have positive integer
valuation for every $n\ge1$. Both points are formal at every place
above $p$, so the already proved two-place proxy lemma applies.
In particular
\[
 v_p(T)>\frac{14M+7}{8}R
       =\left(7\cdot2^{n-1}-\frac78\right)R.
\]
Since $R\ge1$ and $v_p(T)$ is an integer,
\[
 \boxed{\quad v_p(a-1)\ge v_p(T)
              \ge 7\cdot2^{n-1}.\quad}
\]
The first inequality uses
$v_p(T)=v_p((a-1)/h)$ with nonzero integral $h$.
This is a uniform radical exponent; the stronger inequality
$v_p(T)\ge(7\cdot2^{n-1})R$ is not asserted.
If one starts with only oddness of $R$, the constant all-left path
first excludes $R<0$, exactly as in the selected proof.

\paragraph{Completeness.}
For a fixed integral input and $p\mid r$, require
\[
 v_p(A),v_p(B),v_p(C)\ge R,\qquad
 v_p(D_0)\ge uR+1,\qquad v_p(y)\ge3uR+1.
\]
The left subtree satisfies $L\equiv1\pmod p$. On the right,
\[
 v_p(J)\ge2uR+2=\frac{M+1}{2}R+2,
 \qquad v_p(\kappa J^2)\ge R+4>0.
\]
Hence $W\equiv1\pmod p$. The bound on $y$ is simply the finite
requirement $v_p(T)\ge7uR+1$. It is attainable by choosing the scalar
point sufficiently deep after the shared point is fixed. At primes
dividing 42 but not $r$, require the nonconstant leaves to be divisible
by $p$. This again gives $W\equiv1$. Real positivity gives $W\ge1$,
so the same pole realization and continuity argument applies. The
stronger finite preparation does not change the order of prime choices.

The first four values of the uniform order are therefore
\[
\begin{array}{c|rrrr}
 n\text{ Laurent coordinates after }r^{-1}&1&2&3&4\\
 M&3&7&15&31\\
 \text{uniform scalar order}&7&14&28&56
\end{array}
\]
The selected $n=3$ implementation has the separately measured count
of 399 variables in isolation. The lemma itself does not claim final
compiler counts for the other members of this family.
